\section{Introduction}
\label{sec:introduction}

Connected laboratories join instruments, sample movement, scheduling, data capture, and increasingly autonomous decision processes into one operational system. Standards and open frameworks have reduced some integration costs, including SiLA and SiLA~2 for device communication and PyLabRobot for hardware-agnostic programming \cite{bar2012sila,juchli2022sila2,porr2020gateway,bromig2022manager,wierenga2023pylabrobot}. Self-driving laboratory architectures extend this stack through adaptive experiment selection, automated synthesis and characterization, and generalized scheduling \cite{abolhasani2023rise,hung2024community,canty2023autonomy,fei2024alabos,cousty2024glas}. Their practical reliability depends on more than successful command dispatch. Interfaces must be accessible, deployments must be identifiable, physical resources must be represented with validated semantics, execution must remain observable, and interrupted runs must reach a scientifically defensible disposition.

The connected-laboratory literature principally describes standards, architectures, algorithms, and exemplar systems. Operational integration knowledge appears less systematically. Practitioners still need to determine which vendor artifacts constitute deployable source, whether a nominally available driver can be licensed and tested, why a shared labware definition fails on another machine, which scientific constraints a scheduler requires, and what material state remains after a partially executed operation. Public technical discussions preserve these details together with competing explanations, workarounds, vendor clarifications, and cases in which local diagnosis becomes a reusable public artifact.

These discussions are scientifically valuable only under explicit inferential limits. A thread count has no run denominator and cannot measure field-wide incidence. Reporting propensity differs across products, installed bases, user communities, and problem classes. Public discussions can terminate after work moves to vendor support, direct messages, meetings, or local repositories. Product developers can report useful quantitative evidence whose scope stops at simulation or another bounded validation stage. The defensible objective is therefore to identify recurrent engineering constructs, operationalize them over explicit case units, and use the resulting mechanisms to define prospective measurements for connected laboratories.

This paper addresses three research questions.
\begin{enumerate}[label=\textbf{RQ\arabic*.}]
  \item Which recurrent technical and ecosystem bottlenecks are visible across public connected-laboratory discussions?
  \item Which quantitative measures can be computed over bounded forum-derived objects without treating discussion recurrence as operational incidence?
  \item Which observed relationships and community artifacts define the strongest priorities for connected-laboratory engineering and evaluation?
\end{enumerate}

The study makes four contributions. First, it develops a ten-construct taxonomy spanning ecosystem knowledge, interfaces and representations, runtime coordination, evaluation semantics, and AI-specific context. Second, it introduces bounded measures for integration accessibility, deployment identity, physical-resource completeness, observability, preflight preventability, scheduling constraints, test--claim alignment, and AI context expansion. Third, it evaluates the robustness of these measures through alternative partial-score weights, one-thread recoding, leave-one-thread-out association analysis, and alternative denominator definitions. Fourth, it translates the findings into implementable public infrastructure: a minimum reproducible troubleshooting-question schema, a governed resolved-knowledge record, and proposed interface, event, resource, and scheduling registries.

% claim: C11
The central finding is that connected-laboratory failures concentrate at representational boundaries. Public knowledge must become a maintained artifact. A listed driver must become an accessible and testable interface. Method source must become an identified deployment. A geometry value must become a physically validated resource definition. Telemetry must become adjudicable evidence. A partial execution trace must become a verified recovery decision. AI-generated syntax must become an experimentally evaluated method. This boundary-centered structure links otherwise heterogeneous incidents and identifies where connected-laboratory systems require explicit contracts, identifiers, and validation stages.
